
définissant @@M0@@, et soit @@M1@@ le point générique de @@M2@@. Comme @@M3@@ est intègre, @@M4@@ est l'idéal maximal de @@M5@@, et, puisque @@M6@@ est régulier, @@M7@@ peut être engendré par une suite @@M8@@-régulière. On peut donc trouver une suite de sections de @@M9@@ dans un voisinage de @@M10@@, formant une suite @@M11@@-régulière. Par suite, il existe un voisinage ouvert @@M12@@ de @@M13@@, tel que @@M14@@ soit contenu dans @@M15@@, et que @@M16@@ soit un idéal régulier de @@M17@@. L'immersion @@M18@@ est alors une immersion régulière, dont la codimension est @@M19@@ (EGA IV 5.1.2). Par ailleurs, on a la suite exacte (IV\qquad) :
@@M20@@
où @@M21@@ est l'immersion @@M22@@ et @@M23@@ l'immersion @@M24@@ ; on obtient une suite analogue après tensorisation par @@M25@@. Or @@M26@@, et d'après 4.6, l'image @@M27@@ de la classe de @@M28@@ par @@M29@@ est donc dans @@M30@@. Comme @@M31@@ est un @@M32@@-homomorphisme surjectif, il induit une \emph{surjection} @@M33@@, donc @@M34@@ avec @@M35@@ ; il existe donc @@M36@@ tel que :
@@M37@@
Or @@M38@@, et comme @@M39@@, @@M40@@. Par suite, @@M41@@ peut s'exprimer comme combinaison de classes de faisceaux @@M42@@, où les @@M43@@ sont des sous-schémas fermés intègres de @@M44@@, distincts de @@M45@@. D'après l'hypothèse de récurrence, @@M46@@, ce qui achève la démonstration.

\subsection*{BIBLIOGRAPHIE}
\begin{enumerate}[label={[\arabic*]}]
\item Cartan--Eilenberg, Homological Algebra.
\end{enumerate}



% Expose VIII macro additions
\providecommand{\etale}{\'etale}
\providecommand{\Ch}{\operatorname{Ch}}


